% ============================================================
% sec/01introduccion.tex
%
% DEBE CONTENER, sin excepción:
%   1. Contexto y relevancia del problema
%   2. Pregunta de investigación EXPLÍCITA y refutable
%   3. Hipótesis de trabajo
%   4. Objetivos + contribución esperada
%   5. Estructura del documento
%
% La pregunta tiene que poder responderse CON EVIDENCIA al
% terminar el semestre. Una pregunta que no admite refutación
% no sirve.
% ============================================================
\section{Introducción}
\label{sec:introduccion}

\subsection{Contexto y relevancia}
\label{subsec:contexto}
\pc{2--3 párrafos. Cuál es el dominio del problema, por qué importa
técnica y prácticamente, y cuál es el costo concreto de no resolverlo.
Citar literatura del dominio desde aquí, no dejarlo todo para
Trabajos Relacionados.}

\subsection{Planteamiento del problema}
\label{subsec:problema}
\pc{1--2 párrafos. La limitación específica que este trabajo ataca,
acotada desde el dominio general. Es el embudo entre el contexto
amplio y la pregunta puntual.}

\subsection{Pregunta de investigación}
\label{subsec:pregunta}
\begin{quote}
\textbf{PI:} \pc{Una sola pregunta explícita. Debe poder responderse
con evidencia experimental al cerrar el semestre y debe admitir
refutación: tiene que ser posible que la respuesta sea "no".}
\end{quote}

\noindent\textbf{Por qué es refutable.} \pc{Enunciar la condición
concreta bajo la cual la respuesta sería negativa.}

\noindent\textbf{Por qué es medible.} \pc{Qué métrica, sobre qué
partición, con qué unidad de análisis.}

\subsection{Hipótesis de trabajo}
\label{subsec:hipotesis}
\begin{quote}
\textbf{H1:} \pc{Enunciado falsable, con dirección.}
\end{quote}
\begin{quote}
\textbf{H2:} \pc{Segunda hipótesis, si aplica.}
\end{quote}

\noindent\textbf{Criterios de falsación.} \pc{Qué resultado concreto
refutaría o debilitaría cada hipótesis. Este apartado es lo que separa
una hipótesis de una expectativa.}

\subsection{Objetivos}
\label{subsec:objetivos}
\textbf{Objetivo general.} \pc{Una sola oración.}

\noindent\textbf{Objetivos específicos.}
\begin{enumerate}
  \item \pc{Objetivo 1}
  \item \pc{Objetivo 2}
  \item \pc{Objetivo 3}
  \item \pc{Objetivo 4}
  \item \pc{Objetivo 5}
\end{enumerate}

\subsection{Contribución esperada}
\label{subsec:contribucion}
\pc{Qué gana quien lea este trabajo que hoy no existe: una comparación
que nadie ha hecho, una formalización, un protocolo reproducible.
No prometer resultados: prometer aportes.}

\subsection{Estructura del documento}
\label{subsec:estructura}
La Sección~\ref{sec:relacionados} revisa la literatura en tres ejes y
cierra identificando las brechas. La Sección~\ref{sec:baseline} declara
el baseline de literatura y su comparabilidad. La
Sección~\ref{sec:metodologia} presenta la metodología propuesta. La
Sección~\ref{sec:analisistecnico} desarrolla el análisis técnico
específico del track. La Sección~\ref{sec:etica} documenta las
consideraciones éticas preliminares.